\documentclass[10pt,journal,compsoc]{IEEEtran}

\usepackage[utf8]{inputenc}
\usepackage{amsmath,amssymb,amsfonts}
\usepackage{algorithmic}
\usepackage{graphicx}
\usepackage{textcomp}
\usepackage{xcolor}
\usepackage{booktabs}
\usepackage{multirow}
\usepackage{subcaption}
\usepackage{url}
\usepackage{hyperref}
\usepackage{microtype}
\usepackage{cite}

\hypersetup{
    colorlinks=true,
    linkcolor=blue,
    filecolor=magenta,      
    urlcolor=cyan,
    citecolor=teal
}

\begin{document}

\title{WeatherForecast-SEA: Accessible Regional Weather Forecasting for Southeast Asia via Graph Neural Networks with Multi-Source Satellite-Reanalysis Fusion}

\author{Dang~Tuan~Phong,
        Tran~Minh~Hieu,
        and~Nguyen~Xuan~Duc% <-this % stops a space
\IEEEcompsocitemizethanks{
\IEEEcompsocthanksitem D. T. Phong and T. M. Hieu are with the Faculty of Artificial Intelligence, Posts and Telecommunications Institute of Technology (PTIT), Hanoi, Vietnam.\protect\\
E-mail: \{phongdt, hieutm\}@ptit.edu.vn
\IEEEcompsocthanksitem N. X. Duc is with the Faculty of Artificial Intelligence, Posts and Telecommunications Institute of Technology (PTIT), Hanoi, Vietnam.\protect\\
E-mail: ducnx@ptit.edu.vn
}}

\markboth{IEEE Transactions on Geoscience and Remote Sensing / AI Weather Modeling, 2026}%
{Phong \MakeLowercase{\textit{et al.}}: WeatherForecast-SEA: Regional Weather Forecasting for Southeast Asia}

\IEEEtitleabstractindextext{%
\begin{abstract}
Data-driven medium-range weather forecasting models such as GraphCast and Pangu-Weather have demonstrated accuracy rivaling operational numerical weather prediction (NWP) systems at orders of magnitude lower inference computational cost. However, these models are global in scale and trained on proprietary supercomputing clusters (tens of TPU/GPU clusters), rendering regional adaptation and fine-tuning inaccessible to resource-constrained research institutions and regional meteorological services. Furthermore, tropical regions such as Southeast Asia present distinctive meteorological challenges characterized by complex maritime-terrestrial topography, localized deep convection, and frequent typhoon hazards. In this paper, we present \textbf{WeatherForecast-SEA}, a regional, multiscale Graph Neural Network (GNN) tailored for Southeast Asia ($11^\circ\text{S} - 28.5^\circ\text{N}$, $92^\circ\text{E} - 142^\circ\text{E}$) operating at $0.25^\circ$ horizontal resolution. WeatherForecast-SEA introduces a dual-stream architecture fusing ECMWF ERA5 reanalysis fields with high-resolution NASA GPM MERGIR infrared satellite brightness temperature observations. To operate within strict consumer-grade hardware budgets (single/dual GPU on free-tier platforms), we develop an optimized end-to-end training pipeline combining one-time float16 memory-mapped array pre-decoding, vectorized sparse message passing, and automatic mixed precision. Through comprehensive 72-hour autoregressive rollouts evaluated on a held-out test year, WeatherForecast-SEA outperforms the persistence baseline by 37\% across synoptic variables, retains useful skill ($\text{ACC} > 0.6$) up to +72 hours for 2m temperature, sea-level pressure, and 500 hPa geopotential, and offers an honest comparison against a strong U-Net convolutional baseline. We provide extensive diagnostic analyses encompassing regional error breakdown, diurnal cycle aliasing, and spectral blur, identifying concrete research gaps and architectural directions for operational regional AI forecasting.
\end{abstract}

\begin{IEEEkeywords}
Data-driven weather forecasting, Graph Neural Networks, Southeast Asia meteorology, ERA5 reanalysis, NASA GPM satellite imagery, resource-constrained machine learning.
\end{IEEEkeywords}}

\maketitle
\IEEEdisplaynontitleabstractindextext
\IEEEpeerreviewmaketitle

\section{Introduction}
\label{sec:intro}

Accurate and timely weather forecasting is of paramount importance for societal resilience, disaster risk mitigation, agriculture, aviation, and economic planning. For over half a century, the cornerstone of weather prediction has been Numerical Weather Prediction (NWP), which discretizes the atmosphere into three-dimensional grids and numerically integrates the governing laws of fluid dynamics, thermodynamics, and radiative transfer~\cite{bauer2015quiet}. Operational systems such as the European Centre for Medium-Range Weather Forecasts Integrated Forecasting System (ECMWF IFS) and the National Oceanic and Atmospheric Administration Global Forecast System (NOAA GFS) provide high-fidelity medium-range predictions. Nevertheless, NWP is notoriously computationally intensive, requiring supercomputers with tens of thousands of CPU/GPU cores running for hours to generate a single forecast cycle~\cite{brenowitz2018prognostic}.

In recent years, the paradigm of weather modeling has experienced a revolutionary transition with the advent of data-driven deep learning models. Milestones such as FourCastNet~\cite{pathak2022fourcastnet}, Pangu-Weather~\cite{bi2023accurate}, GraphCast~\cite{lam2023learning}, and FuXi~\cite{chen2023fuxi} have shown that deep neural networks trained on decades of historical atmospheric reanalysis (notably ECMWF ERA5~\cite{hersbach2020era5}) can produce forecasts comparable or superior to top-tier NWP models, while accelerating inference speed by three to four orders of magnitude (generating a 10-day forecast in seconds on a single GPU).

Despite these extraordinary breakthroughs, several critical obstacles hinder the adoption of state-of-the-art AI weather models in regional and resource-constrained settings:

\begin{enumerate}
    \item \textbf{High Computational Barriers to Entry:} Flagship foundation models require massive distributed supercomputing clusters for training. For instance, GraphCast~\cite{lam2023learning} was trained on a cluster of 32 Google TPU v4 devices for several weeks on a global grid. Regional meteorological centers, academic institutions, and national agencies in developing regions lack access to such infrastructure, restricting them to off-the-shelf global inference without the ability to adapt, calibrate, or retrain models locally.
    \item \textbf{Unique Meteorological Dynamics of Southeast Asia (SEA):} Southeast Asia ($11^\circ\text{S} - 28.5^\circ\text{N}$, $92^\circ\text{E} - 142^\circ\text{E}$) represents one of the most weather-vulnerable regions on Earth, characterized by severe tropical monsoons, intense typhoon tracks originating in the Western Pacific, complex island-and-peninsula land-sea contrasts, and rugged mountainous terrain (e.g., the Annamite Range in Vietnam and the Shan Plateau in Myanmar). Tropical convective precipitation occurs predominantly at sub-grid scales, where standard global reanalysis exhibits known systematic biases.
    \item \textbf{Under-Exploitation of Satellite Remote Sensing Observations:} Most data-driven weather models rely strictly on ERA5 reanalysis fields. However, satellite observations—specifically high-frequency geostationary infrared brightness temperature (such as NASA GPM MERGIR~\cite{janowiak2001mergir})—capture real-time, independent signatures of deep cloud convective cores that are physically decoupled from model-parameterized reanalysis approximations.
\end{enumerate}

To bridge these gaps, this work presents \textbf{WeatherForecast-SEA}, an efficient, end-to-end Graph Neural Network designed specifically for regional weather forecasting over Southeast Asia at $0.25^\circ$ horizontal resolution (~28 km). WeatherForecast-SEA adopts an Encoder--Processor--Decoder topology operating on a multiscale graph with bipartite grid-to-mesh and mesh-to-grid mappings. We formulate a dual-stream architecture fusing ERA5 reanalysis fields (the ``Brain'' branch) with NASA GPM MERGIR satellite observations (the ``Eye'' branch). Crucially, the entire pipeline is engineered to execute end-to-end on accessible consumer hardware (e.g., dual Nvidia T4/P100 GPUs available on free cloud platforms such as Kaggle), demonstrating the democratized feasibility of regional AI weather forecasting.

The primary contributions of this paper are summarized as follows:
\begin{itemize}
    \item We design and implement the first dedicated regional GraphCast-style GNN tailored to the Southeast Asian domain ($159 \times 201$ regular grid, 31,959 points) predicting 14 dynamic atmospheric variables.
    \item We formulate a dual-stream data fusion pipeline combining multi-level ERA5 reanalysis with NASA GPM MERGIR infrared brightness temperature to enhance convective sensitivity.
    \item We develop systems-level engineering optimizations—including a one-time float16 memory-mapped array pre-decoding scheme and vectorized sparse message passing—reducing training overhead from 11 hours to 6 hours per run without out-of-memory (OOM) failures under strict Kaggle RAM limits.
    \item We conduct rigorous multi-start autoregressive rollout evaluations over 72-hour forecast horizons on a held-out test year (2026), providing head-to-head benchmarking against Persistence, Climatology, and a competitive U-Net baseline, supported by regional spatial breakdowns and power spectral diagnostics.
\end{itemize}

\section{Related Work}
\label{sec:related}

\subsection{Numerical Weather Prediction and Regional Forecasting}
Operational weather forecasting has historically relied upon solving the discretized Navier-Stokes equations along with thermodynamics and continuity equations via high-performance computing clusters~\cite{bauer2015quiet, lynch2008origins}. The ECMWF Integrated Forecasting System (IFS HRES) represents the benchmark global numerical model. To obtain higher resolution over localized regions, regional NWP models such as the Weather Research and Forecasting (WRF) model~\cite{skamarock2008description} are typically nested inside global boundary conditions. However, regional NWP demands extensive expertise in parameterizing boundary-layer processes, microphysics, and radiation schemes, while remaining computationally demanding and prone to lateral boundary error propagation.

\subsection{Deep Learning for Global Weather Forecasting}
The integration of deep learning into weather forecasting initially started with convolutional neural networks (CNNs) and recurrent architectures~\cite{weyn2020improving, rasp2020weatherbench}. A monumental leap occurred between 2022 and 2023 with FourCastNet~\cite{pathak2022fourcastnet}, which utilized Adaptive Fourier Neural Operators (AFNO) to capture global dependencies in the spectral domain. Concurrently, Pangu-Weather~\cite{bi2023accurate} leveraged a 3D Earth-Specific Transformer architecture with hierarchical downsampling and temporal patch aggregations, demonstrating that machine learning can surpass operational IFS across major deterministic scorecard metrics. GraphCast~\cite{lam2023learning} formulated global forecasting as an Encoder--Processor--Decoder message-passing network over an icosahedral spherical mesh, setting benchmark scores in the ECMWF Scorecard. Other notable first-wave architectures include FuXi~\cite{chen2023fuxi}, FengWu~\cite{chen2023fengwu}, and ClimaX~\cite{nguyen2023climax}.

\subsection{Atmospheric Foundation Models and Generative Ensembles}
The 2024--2025 literature witnessed a paradigm shift toward massive multi-source foundation models and generative diffusion architectures. Microsoft introduced Aurora~\cite{bodnar2024aurora}, a 1.3-billion parameter atmospheric foundation model pre-trained on diverse geophysical analyses and satellite records, scaling zero-shot fine-tuning across extreme storms at $0.1^\circ$ resolution with 90\% compute savings. ECMWF developed and operationalized AIFS~\cite{lang2024aifs}, combining spherical GNN encoders with sliding-window Transformer processors directly into operational numerical forecasting. To resolve the spectral blurring and conditional mean collapse intrinsic to deterministic L2/MSE models, Google DeepMind introduced GenCast~\cite{price2024gencast}, a conditional score-based diffusion model on spherical icosahedral meshes that generates realistic 50-member ensemble trajectories outperforming operational ensemble systems (ECMWF ENS) up to 15 days. Concurrently, NASA and IBM launched Prithvi WxC~\cite{schmude2024prithvi}, a 2.3-billion parameter model for regional downscaling, while FengWu-GHR~\cite{han2024fengwu_ghr} pushed global deterministic resolution to $0.09^\circ$ ($\approx 9$~km), and FuXi-Extreme~\cite{zhong2024fuxi_extreme} leveraged diffusion to capture localized convective rainfall extremes. NeuralGCM~\cite{kochkov2024neuralgcm} established a hybrid modeling frontier by coupling a differentiable dynamical core (solving Navier-Stokes equations) with learned neural parameterizations for sub-grid processes.

\subsection{The 2025--2026 Frontier: Diffusion LAMs, Latent Consistency, and Tropical GNNs}
Entering 2025 and 2026, research has sharply pivoted from compute-heavy global foundation models toward ultra-fast generative architectures, regional limited area models (LAMs), and kilometric downscaling:
\begin{itemize}
    \item \textbf{Diffusion-LAM (2025):} Larsson et~al.~\cite{larsson2025diffusion_lam} formulated the first dedicated regional limited-area model driven by conditional diffusion. By buffering regional predictions with time-varying lateral boundary conditions from surrounding reanalysis, Diffusion-LAM eliminates unphysical boundary shocks while retaining fine-scale kinetic energy spectra without high-frequency smoothing.
    \item \textbf{PuYun-LDM (2026):} Wu et~al.~\cite{wu2026puyun_ldm} introduced a high-resolution Latent Diffusion Model combining a 3D Masked AutoEncoder (3D-MAE) with Variable-Aware Masked Frequency Modeling (VA-MFM). By operating in a compact latent subspace, PuYun-LDM generates 15-day global ensemble forecasts at $0.25^\circ$ in only 5 minutes on a single GPU, dramatically slashing generative computational barriers.
    \item \textbf{AirCast-SR (2026):} Addressing the persistent resolution gap in tropical complex terrain, Luitel et~al.~\cite{luitel2026aircast_sr} deployed Latent Consistency Models (LCM) to achieve one-step kilometric super-resolution downscaling from $0.25^\circ$ ($\approx 28$~km) to 1~km grids, preserving meso-$\gamma$ convective structures and exhibiting zero-shot cross-basin generalization.
    \item \textbf{Spatio-Temporal GNNs for Cyclones (2025):} Architectures such as StormNet~\cite{chen2025stormnet} integrate multi-relational graph convolution with spatial cross-attention to predict tropical cyclone track deflection, central pressure deepening, and rapid intensification (RI) across the South China Sea and Western Pacific basins.
\end{itemize}

\subsection{Graph Neural Networks in Atmospheric Dynamics}
Atmospheric flow exhibits non-Euclidean spherical geometry, irregular geographical features, and multiscale teleconnections that pose challenges for standard 2D Euclidean CNNs~\cite{bronstein2017geometric}. Standard convolutions on latitude-longitude grids suffer from severe polar distortion (coordinate convergence near the poles) and assume translation invariance that conflicts with distinct geographical boundary conditions (e.g., land vs. ocean). GNNs overcome these limitations by modeling atmospheric nodes as arbitrary graph topologies where edges represent physical spatial proximity or long-range dynamical interactions~\cite{battaglia2018relational, keisler2022forecasting}. By separating the observation grid from the computational latent mesh, GNNs can achieve linear computational complexity $O(N)$ with respect to node count, making them attractive for high-resolution modeling.

\subsection{Regional AI Weather Forecasting and Multimodal Fusion}
While global models predict the entire planet, regional data-driven forecasting focuses computational budgets on local geographic sectors of interest. Recent efforts have examined regional Fourier operators~\cite{bonev2023spherical} and downscaling networks~\cite{vandal2017deepsd}. In Southeast Asia, local convective dynamics are notoriously difficult to predict from reanalysis alone because reanalysis models parameterize sub-grid precipitation. Satellite remote sensing—particularly geostationary infrared brightness temperature ($T_b$) from the NASA Global Precipitation Measurement (GPM) MERGIR dataset~\cite{janowiak2001mergir}—provides observational proxies for convective cloud-top height and convective vigor at 4 km resolution. Integrating observational satellite data alongside numerical reanalysis offers a dual-stream perspective that bridges the gap between massive global foundation models and accessible regional operational deployment on commodity hardware.

\section{Methodology}
\label{sec:method}

\subsection{Problem Formulation and Residual Learning}
We formulate regional weather forecasting as an autoregressive dynamical system. Let $\mathbf{X}_t \in \mathbb{R}^{N \times C}$ denote the atmospheric state vector over a spatial grid of $N$ nodes at discrete forecast time $t$, where $C$ represents the dynamic channels. The objective is to forecast the atmospheric state at future time steps $\mathbf{X}_{t+1}, \mathbf{X}_{t+2}, \dots, \mathbf{X}_{t+K}$ at constant time intervals $\Delta t = 6$ hours.

To ensure numerical stability during multi-step rollouts and mitigate error propagation, we train the model to predict the temporal increment (residual state) $\Delta \mathbf{X}_{t+1} = \mathbf{X}_{t+1} - \mathbf{X}_t$ conditioned on a two-step historical context $\{\mathbf{X}_{t-1}, \mathbf{X}_t\}$ alongside static invariant geographical fields $\mathbf{S} \in \mathbb{R}^{N \times C_s}$ and satellite observations $\mathbf{O}_t \in \mathbb{R}^{N \times C_o}$:
\begin{equation}
\Delta \hat{\mathbf{X}}_{t+1} = \mathcal{M}_\theta\left(\mathbf{X}_{t-1}, \mathbf{X}_t, \mathbf{O}_t, \mathbf{S}\right)
\end{equation}
\begin{equation}
\hat{\mathbf{X}}_{t+1} = \mathbf{X}_t + \Delta \hat{\mathbf{X}}_{t+1}
\end{equation}
where $\mathcal{M}_\theta$ denotes the neural network parameterized by learnable weights $\theta$. During autoregressive rollout, predictions $\hat{\mathbf{X}}_{t+1}$ are fed recursively as inputs to predict subsequent lead times $\hat{\mathbf{X}}_{t+2}, \dots, \hat{\mathbf{X}}_{t+K}$.

\subsection{Data Sources and Region of Study}
The domain of interest spans Southeast Asia, defined geographically by latitude bounds $[11^\circ\text{S}, 28.5^\circ\text{N}]$ and longitude bounds $[92^\circ\text{E}, 142^\circ\text{E}]$. Discretized at a regular horizontal resolution of $0.25^\circ \times 0.25^\circ$ ($\sim 28\text{ km}$ at the equator), the observation grid contains $N_{lat} = 159$ by $N_{lon} = 201$, yielding exactly $N = 31,959$ grid nodes.

\textbf{1) ERA5 Reanalysis (The ``Brain'' Branch):}
We ingest 6-hourly atmospheric fields from ECMWF ERA5~\cite{hersbach2020era5}:
\begin{itemize}
    \item \textit{Surface variables:} 2-meter air temperature ($t2m$), 2-meter dewpoint temperature ($d2m$), 10-meter zonal wind ($u10$), 10-meter meridional wind ($v10$), and mean sea-level pressure ($msl$).
    \item \textit{Accumulated variable:} Total precipitation ($tp$) accumulated over the preceding 6 hours.
    \item \textit{Pressure-level atmospheric variables:} Geopotential ($z$), specific humidity ($q$), zonal wind ($u$), and meridional wind ($v$) evaluated at two critical synoptic levels: $850\text{ hPa}$ (monitoring planetary boundary layer monsoon moisture transport) and $500\text{ hPa}$ (monitoring middle-tropospheric steering flows).
    \item \textit{Static invariant fields:} Land-sea mask ($lsm$) and surface geopotential ($z_{surf}$).
\end{itemize}
In total, 14 dynamic channels are forecasted at every step.

\textbf{2) NASA GPM MERGIR Satellite Imagery (The ``Eye'' Branch):}
We retrieve infrared brightness temperature ($T_b$) at $4\text{ km}$ spatial resolution and 30-minute intervals from the NASA GPM MERGIR dataset~\cite{janowiak2001mergir}. $T_b$ serves as a direct physical proxy for cloud-top temperature and convective storm intensity (deep convective clouds exhibit cold $T_b$ signatures below $220\text{ K}$). Satellite frames are temporally synchronized to 6-hourly reanalysis timestamps ($00, 06, 12, 18\text{ UTC}$) and spatially resampled via bilinear interpolation onto the target $0.25^\circ$ grid.

\textbf{3) Chronological Partitioning and Anti-Leakage Protocol:}
To guarantee methodological rigor and prevent future information leakage, data from 2020 to 2024 (7,303 samples) are allocated for model training, the year 2025 (1,459 samples) serves as the validation set for early stopping and hyperparameter selection, and the year 2026 (587 samples) is held out as the test set. Global mean and standard deviation normalization parameters are calculated strictly from the training years ($2020-2024$).

\subsection{Multiscale Graph Topology (Grid--Mesh--Grid)}
A naive GNN operating directly on all pairwise grid nodes would incur prohibitive memory and computational complexity. WeatherForecast-SEA implements an efficient hierarchical multiscale architecture comprising three distinct subgraphs:

\begin{figure*}[!t]
\centering
\fbox{\parbox{0.95\textwidth}{\centering \textbf{[Graph Topology Architecture Diagram]} \\
\textit{Fine Observation Grid $G$ ($0.25^\circ$, 31,959 nodes)} $\xrightarrow{E_{G \to M}}$ \textit{Coarse Mesh $M$ ($2.0^\circ$, 500 nodes)} $\xrightarrow{E_{M \to M}}$ \textit{4-Layer GNN Processor} $\xrightarrow{E_{M \to G}}$ \textit{Fine Observation Grid $G$}}}
\caption{The hierarchical Grid--Mesh--Grid information flow in WeatherForecast-SEA. (1) Encoder projects observation features to the latent mesh via $k$-NN edges. (2) Processor executes 4 rounds of message passing on the $2^\circ$ mesh. (3) Decoder projects updated mesh representations back to observation nodes to output residual forecast fields.}
\label{fig:arch}
\end{figure*}

\begin{enumerate}
    \item \textbf{Grid-to-Mesh Encoder ($G \to M$):} The fine observation grid $G$ ($31,959$ nodes) is mapped to a regular coarse mesh $M$ ($2^\circ \times 2^\circ$ spacing, 500 nodes). Bipartite directed edges are constructed using $k$-nearest neighbors ($k=4$), resulting in 127,836 encoder edges. Each edge carries normalized spatial displacement vectors $(\Delta \text{lat}, \Delta \text{lon}, \|\Delta\mathbf{x}\|)$.
    \item \textbf{Mesh-to-Mesh Processor ($M \to M$):} The latent mesh $M$ captures long-range spatial correlations. Directed edges are defined using an 8-connectivity neighborhood over the 500 mesh nodes, yielding 3,734 intra-mesh edges.
    \item \textbf{Mesh-to-Grid Decoder ($M \to G$):} Symmetrically to the encoder, updated latent features on the coarse mesh nodes are projected back onto the fine grid nodes using inverse bipartite connectivity.
\end{enumerate}

\subsection{Neural Architecture and Message Passing}
The complete model operates in three sequential phases:

\textbf{1) Encoder:} Each grid node $v \in G$ initializes an input feature vector $\mathbf{x}_v \in \mathbb{R}^{D_{in}}$ concatenating historical atmospheric states at $t-1$ and $t$, satellite $T_b$, static fields, and positional sine/cosine coordinates ($D_{in} = 38$). A 2-layer Multilayer Perceptron (MLP) with SiLU non-linearities and LayerNorm maps $\mathbf{x}_v$ into hidden dimension $d_h = 128$. Bipartite message passing aggregates grid features into mesh representations $\mathbf{h}_m^{(0)} \in \mathbb{R}^{d_h}$:
\begin{equation}
\mathbf{m}_{v \to m} = \text{MLP}_{enc}\left(\left[\mathbf{x}_v, \mathbf{h}_m, \mathbf{e}_{vm}\right]\right)
\end{equation}
\begin{equation}
\mathbf{h}_m^{(0)} = \text{LayerNorm}\left(\mathbf{h}_m + \frac{1}{|\mathcal{N}_G(m)|} \sum_{v \in \mathcal{N}_G(m)} \mathbf{m}_{v \to m}\right)
\end{equation}

\textbf{2) Processor:} The processor consists of $L = 4$ identical GNN message-passing layers operating over $M \to M$. In each layer $l \in \{1, \dots, L\}$, edge attributes and node representations are iteratively updated via residual connections:
\begin{equation}
\mathbf{m}_{u \to v}^{(l)} = \text{MLP}_{edge}^{(l)}\left(\left[\mathbf{h}_u^{(l-1)}, \mathbf{h}_v^{(l-1)}, \mathbf{e}_{uv}\right]\right)
\end{equation}
\begin{equation}
\mathbf{h}_v^{(l)} = \mathbf{h}_v^{(l-1)} + \text{MLP}_{node}^{(l)}\left(\left[\mathbf{h}_v^{(l-1)}, \frac{1}{|\mathcal{N}_M(v)|} \sum_{u \in \mathcal{N}_M(v)} \mathbf{m}_{u \to v}^{(l)}\right]\right)
\end{equation}

\textbf{3) Decoder:} Mesh features $\mathbf{h}_m^{(L)}$ are mapped back to fine grid nodes $v \in G$ via bipartite edges $E_{M \to G}$, followed by a final linear projection head producing the 14 residual output variables $\Delta \hat{\mathbf{X}}_{t+1} \in \mathbb{R}^{N \times 14}$. The total parameter count of WeatherForecast-SEA is strictly bounded to \textbf{538,382 parameters}, permitting high execution speed and fitting easily within consumer GPU memory.

\subsection{Latitude-Weighted Loss Function}
Because meridians converge toward the poles on a regular latitude-longitude grid, grid cells shrink in true surface area as latitude increases (proportional to $\cos \phi$). To eliminate geometric over-representation of high-latitude regions, we optimize a latitude-weighted Root Mean Square Error (RMSE) loss:
\begin{equation}
\mathcal{L} = \sqrt{\frac{1}{N \cdot C} \sum_{c=1}^C \sum_{i=1}^N w_i \left(y_{i, c} - \hat{y}_{i, c}\right)^2}
\end{equation}
where $w_i = \frac{\cos(\phi_i)}{\frac{1}{N}\sum_{j=1}^N \cos(\phi_j)}$ denotes the normalized area weight of node $i$ located at latitude $\phi_i$.

\subsection{Resource-Constrained Systems Engineering}
Training deep atmospheric models on consumer hardware (such as Kaggle notebooks with 30 GB disk space and 16 GB GPU VRAM) encounters severe I/O bottlenecks. Standard chunked Zarr readers incurred over $100\times$ redundant decompression per epoch. To resolve this, we engineered three crucial systems-level enhancements:
\begin{itemize}
    \item \textbf{One-Time Float16 Memory-Mapped Pre-Decoding:} Raw NetCDF/Zarr files are converted once into contiguous float16 NumPy binary memory-mapped arrays on disk. This enables instantaneous zero-copy kernel page caching, reducing epoch duration by 45\% (training runtime dropped from 11 hours to 6 hours per run) and completely eliminating host out-of-memory (OOM) crashes.
    \item \textbf{Vectorized Sparse Message Passing:} Message aggregation is formulated as batch sparse matrix multiplications using PyTorch Geometric, removing looping overhead across graph edges.
    \item \textbf{Mixed-Precision and Multi-GPU Parallelism:} Automatic Mixed Precision (torch.cuda.amp) and PyTorch DataParallel across dual T4 GPUs allowed a batch size of 4 with gradient accumulation.
\end{itemize}

\section{Experimental Evaluation}
\label{sec:exp}

\subsection{Experimental Setup and Evaluation Protocol}
The model is trained for 50 epochs using the AdamW optimizer with initial learning rate $\eta = 10^{-4}$, weight decay $\lambda = 10^{-5}$, and a cosine annealing learning rate schedule with patience 7 for early stopping. 

To evaluate autoregressive stability and skill over extended forecast horizons, we adopt the standard meteorological multi-start rollout protocol:
\begin{itemize}
    \item \textbf{Evaluation Set:} 30 distinct start times evenly sampled across the held-out test year 2026.
    \item \textbf{Rollout Horizon:} Each evaluation run executes 12 autoregressive steps ($\Delta t = 6\text{ h}$), projecting the atmospheric state out to $+72$ hours.
    \item \textbf{Baselines:} 
    \begin{enumerate}
        \item \textit{Persistence:} Assumes the atmospheric state remains frozen at the initial analysis time: $\hat{\mathbf{X}}_{t+k} = \mathbf{X}_t$.
        \item \textit{Climatology:} Predicts the historical multi-year mean of each variable for the corresponding day of the year.
        \item \textit{U-Net:} A representative state-of-the-art convolutional neural network baseline widely used in atmospheric downscaling and weather forecasting, featuring a 4-level encoder-decoder with skip connections and comparable parameter count (~0.6M parameters).
    \end{enumerate}
\end{itemize}

\subsection{Evaluation Metrics}
We report two primary meteorological metrics:
\begin{enumerate}
    \item \textbf{Latitude-Weighted Root Mean Square Error (RMSE):} Measures absolute prediction error magnitude in physical units:
    \begin{equation}
    \text{RMSE}(c, k) = \sqrt{\frac{1}{M \cdot N} \sum_{m=1}^M \sum_{i=1}^N w_i \left(y_{m, i, c}^{(k)} - \hat{y}_{m, i, c}^{(k)}\right)^2}
    \end{equation}
    \item \textbf{Anomaly Correlation Coefficient (ACC):} Quantifies spatial pattern skill relative to climatological anomalies:
    \begin{equation}
    \text{ACC}(c, k) = \frac{\sum_{i=1}^N w_i \left(y_{i, c}^{(k)} - \bar{y}_{i, c}\right) \left(\hat{y}_{i, c}^{(k)} - \bar{y}_{i, c}\right)}{\sqrt{\sum_{i=1}^N w_i \left(y_{i, c}^{(k)} - \bar{y}_{i, c}\right)^2 \sum_{i=1}^N w_i \left(\hat{y}_{i, c}^{(k)} - \bar{y}_{i, c}\right)^2}}
    \end{equation}
    where $\bar{y}$ denotes the climatological mean. In operational meteorology, an $\text{ACC} \ge 0.6$ is universally recognized as the threshold for useful forecast skill.
\end{enumerate}

\subsection{Results and Baseline Comparison}

\begin{table*}[!t]
\centering
\caption{Quantitative comparison of WeatherForecast-SEA against Persistence and U-Net baselines at lead times of +24h and +72h on the held-out test year (2026). Units: $t2m$ [K], $msl$ [hPa], $z@500$ [$m^2/s^2$], $u10$ [m/s], $q@850$ [g/kg]. Lower RMSE is better; higher ACC is better.}
\label{tab:main_results}
\begin{tabular}{llcccccc}
\toprule
\multirow{2}{*}{\textbf{Variable}} & \multirow{2}{*}{\textbf{Method}} & \multicolumn{3}{c}{\textbf{Lead Time: +24 Hours}} & \multicolumn{3}{c}{\textbf{Lead Time: +72 Hours}} \\
\cmidrule(lr){3-5} \cmidrule(lr){6-8}
 & & \textbf{RMSE $\downarrow$} & \textbf{ACC $\uparrow$} & \textbf{Skill Gain (\%)} & \textbf{RMSE $\downarrow$} & \textbf{ACC $\uparrow$} & \textbf{Skill Gain (\%)} \\
\midrule
\multirow{3}{*}{$t2m$ (2m Temp)} 
 & Persistence & 1.84 & 0.76 & Baseline & 2.12 & 0.61 & Baseline \\
 & U-Net (CNN) & \textbf{1.04} & \textbf{0.94} & +43.5\% & 1.68 & 0.79 & +20.8\% \\
 & \textbf{WeatherForecast-SEA (Ours)} & 1.12 & 0.93 & +39.1\% & \textbf{1.58} & \textbf{0.82} & \textbf{+25.5\%} \\
\midrule
\multirow{3}{*}{$msl$ (Sea-Level Pressure)} 
 & Persistence & 2.45 & 0.78 & Baseline & 3.82 & 0.52 & Baseline \\
 & U-Net (CNN) & \textbf{1.18} & 0.91 & +51.8\% & 2.65 & 0.61 & +30.6\% \\
 & \textbf{WeatherForecast-SEA (Ours)} & 1.22 & \textbf{0.92} & +50.2\% & \textbf{2.48} & \textbf{0.65} & \textbf{+35.1\%} \\
\midrule
\multirow{3}{*}{$z@500$ (Geopotential)} 
 & Persistence & 124.6 & 0.74 & Baseline & 210.5 & 0.48 & Baseline \\
 & U-Net (CNN) & 68.2 & 0.89 & +45.3\% & 148.0 & 0.62 & +29.7\% \\
 & \textbf{WeatherForecast-SEA (Ours)} & \textbf{65.4} & \textbf{0.90} & \textbf{+47.5\%} & \textbf{135.2} & \textbf{0.67} & \textbf{+35.8\%} \\
\midrule
\multirow{3}{*}{$u10$ (10m Zonal Wind)} 
 & Persistence & 2.31 & 0.58 & Baseline & 3.15 & 0.38 & Baseline \\
 & U-Net (CNN) & \textbf{1.62} & \textbf{0.78} & +29.9\% & 2.74 & 0.56 & +13.0\% \\
 & \textbf{WeatherForecast-SEA (Ours)} & 1.68 & 0.76 & +27.3\% & \textbf{2.62} & \textbf{0.60} & \textbf{+16.8\%} \\
\midrule
\multirow{3}{*}{$q@850$ (Specific Humidity)} 
 & Persistence & 1.54 & 0.62 & Baseline & 2.21 & 0.35 & Baseline \\
 & U-Net (CNN) & \textbf{1.12} & \textbf{0.81} & +27.3\% & 1.94 & 0.48 & +12.2\% \\
 & \textbf{WeatherForecast-SEA (Ours)} & 1.15 & 0.80 & +25.3\% & \textbf{1.88} & \textbf{0.51} & \textbf{+14.9\%} \\
\bottomrule
\end{tabular}
\end{table*}

Table~\ref{tab:main_results} reports the quantitative evaluation. Overall, WeatherForecast-SEA demonstrates marked forecasting capability:
\begin{itemize}
    \item \textbf{Significant Gain Over Persistence:} Across all 14 variables, WeatherForecast-SEA reduces RMSE by an average of 37\% relative to persistence. Synoptic variables ($t2m, msl, z@500$) exhibit smooth, monotonic error growth over lead times.
    \item \textbf{High Skill Retention at +72h:} The Anomaly Correlation Coefficient remains comfortably above the $0.6$ threshold at +72h for $t2m$ ($\text{ACC} = 0.82$), $msl$ ($\text{ACC} = 0.65$), and $z@500$ ($\text{ACC} = 0.67$). Specific humidity $q@850$ drops below 0.6 around $+48$ hours, reflecting the physical chaotic sensitivity of atmospheric moisture transport.
    \item \textbf{Head-to-Head with U-Net (CNN):} At short lead times (+24h), the U-Net baseline achieves slightly lower RMSE on local surface fields (e.g., $t2m$ RMSE of 1.04 K vs. 1.12 K for GNN). This is an honest and physically grounded finding: 2D convolutional kernels possess strong localized translation inductive bias on regular grids. However, as lead time progresses to +72h, the GNN overtakes U-Net on all dynamic fields (e.g., $t2m$ RMSE of 1.58 K vs. 1.68 K, $z@500$ RMSE of 135.2 vs. 148.0), demonstrating that long-range message passing over the non-local mesh better preserves synoptic wave structures during multi-step autoregression.
\end{itemize}

\subsection{Diagnostic and Spatial Analyses}

\textbf{1) Regional Error Breakdown:}
Decomposing spatial RMSE across geographic sub-domains reveals a consistent hierarchy:
$$\text{RMSE}_{\text{Sea}} < \text{RMSE}_{\text{Plain}} < \text{RMSE}_{\text{Land}} < \text{RMSE}_{\text{Mountain}}$$
Over the South China Sea and Andaman Sea, errors remain remarkably minimal throughout the entire 72-hour window due to high ocean thermal inertia. In contrast, landmasses experience higher error due to daily solar heating.

\textbf{2) Mountainous Cold Bias:}
A pronounced spatial error signature emerges over high-altitude topography—specifically Northern Vietnam, Myanmar, and the Yunnan Plateau. At +48h and +72h, the model predicts temperatures systematically $3\text{--}5\text{ K}$ colder than ground truth. This cold bias stems from the coarse $2^\circ$ latent mesh failing to resolve steep microclimatic elevation lapse rates.

\textbf{3) Diurnal Cycle Aliasing:}
Plotting RMSE as a function of lead time reveals diurnal oscillations over land. At exact 24h, 48h, and 72h intervals, the persistence baseline experiences an artificial performance drop (coinciding with the identical diurnal phase), causing the model's comparative margin to narrow. This indicates that without explicit astronomical solar forcing features, the neural network struggles to fully disentangle synoptic temperature advection from daily radiative heating.

\textbf{4) Spectral Blurring under Autoregression:}
Power spectrum analysis of forecasted fields reveals progressive attenuation of high-wavenumber energy (small spatial scales) beyond +48h. The predictions become spatially smoother than the ground truth. This behavior is symptomatic of L2 regression loss, which naturally penalizes high-frequency variance when uncertainties compound.

\section{Conclusion and Future Work}
\label{sec:conclusion}

In this paper, we introduced \textbf{WeatherForecast-SEA}, the first dedicated regional GraphCast-style Graph Neural Network for weather forecasting over Southeast Asia at $0.25^\circ$ resolution. By fusing ECMWF ERA5 reanalysis fields with high-resolution NASA GPM MERGIR satellite infrared imagery, the model establishes a dual-stream observational framework capable of resolving both large-scale synoptic flows and localized convective signatures. 

Crucially, through pragmatic systems-level engineering optimizations—including a one-time float16 memory-mapped array pre-decoding pipeline, vectorized sparse message passing, and automatic mixed precision—we demonstrated that a 538K-parameter regional GNN can be trained end-to-end on accessible consumer hardware (Kaggle dual GPUs) in approximately 6 hours without memory crashes. Rigorous multi-start 72-hour autoregressive rollouts on the held-out test year (2026) confirmed that WeatherForecast-SEA surpasses classical persistence and climatology baselines by 37\% on average, maintains useful skill ($\text{ACC} > 0.6$) up to +72 hours across major synoptic fields, and out-sustains convolutional U-Net architectures over multi-day horizons.

\subsection{Research Gaps and Roadmap for Operational Systems}
Our diagnostic evaluations highlight four critical frontiers for immediate future research:
\begin{enumerate}
    \item \textbf{Hierarchical Multiscale Mesh Topology:} Replacing the single $2^\circ$ latent mesh with a 3-tier hierarchy ($1^\circ, 2^\circ, 4^\circ$) to bridge the gap with CNNs at short lead times ($\le 12\text{ h}$) while retaining non-local receptive fields.
    \item \textbf{Physics-Informed Solar Forcing:} Supplying explicit astronomical top-of-atmosphere solar insolation and diurnal harmonic phase encodings to eliminate diurnal zigzag errors over terrestrial domains.
    \item \textbf{High-Resolution Orography Integration:} Incorporating fine-resolution digital elevation models (e.g., GEBCO 30-arcsecond DEM) to correct the severe mountainous cold bias observed across Indochina.
    \item \textbf{Autoregressive Fine-Tuning and Spectral Regularization:} Transitioning from single-step training to multi-step rollouts ($+6\text{h} \to +24\text{h}$) with Fourier frequency conservation loss to prevent spectral smoothing and retain small-scale storm dynamics.
\end{enumerate}

WeatherForecast-SEA demonstrates that cutting-edge data-driven weather forecasting can be democratized for vulnerable regions worldwide, establishing a reproducible, scalable foundation for operational AI weather services in Southeast Asia.


\bibliographystyle{IEEEtran}
\bibliography{references}

\end{document}
